\begin{frontmatter}

\title{Co-firing copper in a cordierite-type glass-ceramic printed by multi-material vat photopolymerisation: a model-based process design}

% Author list: complete before submission.
\author[aff1]{[First Author]\corref{cor1}}
\ead{[email address]}
\author[aff1]{[Second Author]}
\cortext[cor1]{Corresponding author.}
\affiliation[aff1]{organization={[Department, Institution]},
            city={[City]},
            country={[Country]}}

\begin{abstract}
Multi-material vat photopolymerisation can print a glass-ceramic body and its copper conductors as one green part, but whether such a part can be fired is unknown. Carbon must leave the glass before its pores close, copper must stay metallic, and both materials must shrink compatibly, all in one temperature range. A verified model couples char pyrolysis and steam gasification, viscous sintering and crystallisation of the glass, diffusional sintering of copper, and co-firing mechanics in one and three dimensions. Inputs come from the literature, and the glass was calibrated against the published IBM glass-ceramic/copper process. Of the two atmospheres compared, only steam with hydrogen held at a fixed fraction of the Cu/Cu$_2$O equilibrium, as in the IBM process, removes the char before the glass seals. Fine copper then densifies during the burnout, well before the glass, and the free-strain mismatch reaches 15.6\%. A paste and programme optimised together use an assumed retarding effect of char on copper to release both materials at once, cutting the mismatch to 3.7\%, but this design is fragile. A design that tolerates a factor of two in gasification rate and 15~K in the glass transition holds the mismatch at 9--11\% and the copper at 89\% IACS in a 16~h cycle. In a printed system-in-package with a die cavity, vias and an embedded solenoid, the two designs cut surface distortion from 516 to 114--266~\textmu m; copper plasticity keeps the cooling bow small. Gasification kinetics, glass viscosity, copper sintering viscosity and char yield should be measured first.
\end{abstract}

\begin{keyword}
Co-firing \sep Glass-ceramic \sep Copper metallisation \sep Vat photopolymerisation \sep Binder removal \sep Sintering mismatch
\end{keyword}

\end{frontmatter}
